\documentclass[12pt,a4paper]{article}
\usepackage[T2A]{fontenc}
\usepackage[utf8]{inputenc}
\usepackage[russian]{babel}
\usepackage{amsmath,amssymb}
\usepackage[margin=2cm]{geometry}
\usepackage{microtype}
\setlength{\parindent}{0pt}
\setlength{\parskip}{6pt}
\allowdisplaybreaks
\newcommand{\R}{\mathbb R}
\newcommand{\T}{^{\mathsf T}}
\newcommand{\norm}[1]{\left\lVert#1\right\rVert_2}
\title{Домашнее задание 9}
\author{}
\date{}
\begin{document}
\maketitle

\section*{Задача 1}
\[
f(x)=x_1^2+\frac{(x_1-1)^4}{12}+x_2^2+x_1x_2.
\]
Градиент и матрица Гессе:
\[
g(x)=\begin{pmatrix}2x_1+\dfrac{(x_1-1)^3}{3}+x_2\\[4pt]2x_2+x_1\end{pmatrix},\qquad
G(x)=\begin{pmatrix}2+(x_1-1)^2&1\\1&2\end{pmatrix}.
\]
Матрица $G(x)$ симметрична. Главные угловые миноры:
\[
\Delta_1=2+(x_1-1)^2>0,\qquad
\Delta_2=2\bigl(2+(x_1-1)^2\bigr)-1=3+2(x_1-1)^2>0.
\]
По критерию Сильвестра $G(x)\succ0$ при любом $x\in\R^2$.

\textbf{Шаг из $x^0=(2,8)\T$.} В начальной точке
\[
g(x^0)=\begin{pmatrix}37/3\\18\end{pmatrix},\qquad
G(x^0)=\begin{pmatrix}3&1\\1&2\end{pmatrix}.
\]
Решим $G(x^0)p_0=-g(x^0)$:
\[
\begin{cases}3p_{0,1}+p_{0,2}=-37/3,\\p_{0,1}+2p_{0,2}=-18.\end{cases}
\]
Удвоим первое уравнение и вычтем второе:
\[
5p_{0,1}=-\frac{74}{3}+18=-\frac{20}{3},\qquad
p_{0,1}=-\frac43,\qquad p_{0,2}=-\frac{25}{3}.
\]
Таким образом,
\[
x^1=\begin{pmatrix}2\\8\end{pmatrix}+\begin{pmatrix}-4/3\\-25/3\end{pmatrix}
=\begin{pmatrix}2/3\\-1/3\end{pmatrix}.
\]


\section*{Задача 2}
Для функции Розенброка
\[
f(x_1,x_2)=100(x_2-x_1^2)^2+(1-x_1)^2
\]
оба слагаемых неотрицательны и равны нулю одновременно только при $x^*=(1,1)\T$. Это единственный глобальный минимум.

\textbf{Производные:}
\[
g(x)=\begin{pmatrix}400x_1(x_1^2-x_2)+2(x_1-1)\\200(x_2-x_1^2)\end{pmatrix},
\]
\[
G(x)=\begin{pmatrix}1200x_1^2-400x_2+2&-400x_1\\-400x_1&200\end{pmatrix}.
\]
В каждой точке решаем $Gp=-g$ и сравниваем $p$ с $d=x^*-x$.

\textbf{1. Точка $x=(-6/5,1)\T=(-1.2,1)\T$.}
\[
g=\begin{pmatrix}-1078/5\\-88\end{pmatrix},\quad
G=\begin{pmatrix}1330&480\\480&200\end{pmatrix},\quad\det G=35600.
\]
Из второго уравнения $480p_1+200p_2=88$ получаем $p_2=11/25-(12/5)p_1$.
Подстановка в первое даёт $178p_1=22/5$, откуда
\[
p=\begin{pmatrix}11/445\\847/2225\end{pmatrix}
=\frac1{2225}\begin{pmatrix}55\\847\end{pmatrix},\qquad d=\begin{pmatrix}11/5\\0\end{pmatrix}.
\]
Следовательно,
\[
\cos(p,d)=\frac{55}{\sqrt{55^2+847^2}}\approx0.064799.
\]
\textbf{2. Точка $x=(0,0)\T$.}
\[
g=\begin{pmatrix}-2\\0\end{pmatrix},\quad
G=\begin{pmatrix}2&0\\0&200\end{pmatrix},\quad
\begin{cases}2p_1=2,\\200p_2=0.\end{cases}
\]
Поэтому
\[
p=\begin{pmatrix}1\\0\end{pmatrix},\qquad d=\begin{pmatrix}1\\1\end{pmatrix},\qquad
\cos(p,d)=\frac1{\sqrt2}\approx0.707107.
\]
\textbf{3. Точка $x=(3/2,1)\T=(1.5,1)\T$.}
\[
g=\begin{pmatrix}751\\-250\end{pmatrix},\quad
G=\begin{pmatrix}2302&-600\\-600&200\end{pmatrix},\quad\det G=100400.
\]
Из $-600p_1+200p_2=250$ следует $p_2=5/4+3p_1$.
Тогда $2302p_1-600(5/4+3p_1)=-751$, то есть $502p_1=-1$. Значит,
\[
p=\frac1{1004}\begin{pmatrix}-2\\1249\end{pmatrix},\qquad d=\begin{pmatrix}-1/2\\0\end{pmatrix},\qquad
\cos(p,d)=\frac2{\sqrt{4+1249^2}}\approx0.001601.
\]
Наибольшая косинусная близость получается в $(0,0)\T$, наименьшая --- в $(1.5,1)\T$. Направление Ньютона не обязано совпадать с направлением на минимум.

\section*{Задача 3}
Возьмём функцию и начальную точку
\[
f(x_1,x_2)=\sqrt{1+x_1^2}+\sqrt{1+x_2^2},\qquad x^0=(2,2)\T.
\]
Функция зависит от обеих координат и принадлежит $C^2(\R^2)$, так как подкоренные выражения положительны. Кроме того, $f(x)\ge2=f(0,0)$, поэтому она ограничена снизу.

Для $\phi(t)=\sqrt{1+t^2}$ имеем
\[
\phi'(t)=\frac{t}{\sqrt{1+t^2}},\qquad
\phi''(t)=\frac1{(1+t^2)^{3/2}}.
\]
Следовательно,
\[
g(x)=\begin{pmatrix}\dfrac{x_1}{\sqrt{1+x_1^2}}\\[5pt]\dfrac{x_2}{\sqrt{1+x_2^2}}\end{pmatrix},\qquad
G(x)=\begin{pmatrix}(1+x_1^2)^{-3/2}&0\\0&(1+x_2^2)^{-3/2}\end{pmatrix}.
\]
Все диагональные элементы положительны и не превосходят $1$, поэтому
\[
0\prec G(x)\preceq I\qquad\text{для всех }x\in\R^2.
\]
Следовательно, $f$ выпукла.
Кроме того, $\|G(x)\|_{\mathrm{op}}\le1$. Для любых $x,y\in\R^2$ по основной теореме анализа
\[
g(x)-g(y)=\int_0^1G\bigl(y+t(x-y)\bigr)(x-y)\,dt,
\]
поэтому
\[
\norm{g(x)-g(y)}\le\int_0^1\|G(y+t(x-y))\|_{\mathrm{op}}\norm{x-y}\,dt
\le\norm{x-y}.
\]
Таким образом, градиент глобально липшицев с константой $L=1$.

\textbf{Явная итерация.} Гессиан невырожден в любой конечной точке. Поскольку он диагонален, для $i=1,2$
\[
p_{k,i}=-\frac{x_i^k/\sqrt{1+(x_i^k)^2}}{(1+(x_i^k)^2)^{-3/2}}
=-x_i^k\bigl(1+(x_i^k)^2\bigr).
\]
При полном шаге получаем
\[
x_i^{k+1}=x_i^k+p_{k,i}=-(x_i^k)^3,\qquad i=1,2.
\]
\textbf{Доказательство расходимости.} Индукцией из $x_i^0=2$ следует
\[
x_i^k=(-1)^k2^{3^k}.
\]
Действительно, формула верна при $k=0$, а возведение в куб с последующей сменой знака даёт ту же формулу для $k+1$. Поэтому
\[
\norm{x^k}=\sqrt2\,2^{3^k}\longrightarrow\infty,\qquad
f(x^k)=2\sqrt{1+2^{2\cdot3^k}}\longrightarrow\infty.
\]
Полный метод Ньютона расходится.

\section*{Задача 4}
\[
f(x)=e^{x_1+x_2}+x_1^2+x_2^2.
\]
\textbf{Градиент и гессиан.} Обозначим $s=x_1+x_2$. Тогда
\[
g(x)=\begin{pmatrix}e^s+2x_1\\e^s+2x_2\end{pmatrix},\qquad
G(x)=\begin{pmatrix}e^s+2&e^s\\e^s&e^s+2\end{pmatrix}.
\]
\textbf{Оценка $G(x)\succeq2I$.} Для $u=(1,1)\T$ имеем
\[
G(x)=2I+e^suu\T.
\]
Следовательно, при любом $v\in\R^2$
\[
v\T(G(x)-2I)v=e^s(u\T v)^2=e^s(v_1+v_2)^2\ge0.
\]
Таким образом, $G(x)\succeq2I$; в частности, гессиан всюду положительно определён, а функция $2$-сильно выпукла.

\textbf{Ньютоновский шаг.} При $a_k=e^{x_1^k+x_2^k}$ решаем СЛАУ
\[
\begin{pmatrix}a_k+2&a_k\\a_k&a_k+2\end{pmatrix}
\begin{pmatrix}p_{k,1}\\p_{k,2}\end{pmatrix}
=-\begin{pmatrix}a_k+2x_1^k\\a_k+2x_2^k\end{pmatrix},\qquad
x^{k+1}=x^k+p_k.
\]
Определитель матрицы равен $4(a_k+1)>0$, поэтому система имеет единственное решение на каждом шаге.

\textbf{Достаточна ли положительная определённость для глобальной сходимости?}
Нет. Пример из задачи 3 имеет $G(x)\succ0$ во всех точках, однако полные шаги из $(2,2)\T$ расходятся. Положительная определённость гарантирует направление спуска при $g(x)\ne0$:
\[
g(x)\T p=-g(x)\T G(x)^{-1}g(x)<0,
\]
но не гарантирует уменьшение функции при длине шага $\alpha=1$.

Для данной функции сходимость нужно исследовать отдельно: контрпример из задачи 3 опровергает общее утверждение, но не доказывает расходимость для $e^{x_1+x_2}+x_1^2+x_2^2$.

\textbf{Локальная квадратичная сходимость.} Достаточно, чтобы в окрестности стационарной точки $x^*$ гессиан был липшицев,
\[
\|G(x)-G(y)\|_{\mathrm{op}}\le M\norm{x-y},
\]
и $G(x^*)\succ0$, а начальная точка была достаточно близка к $x^*$. Тогда точные полные шаги корректно определены и
\[
\norm{x^{k+1}-x^*}\le C\norm{x^k-x^*}^2.
\]


\textbf{Зачем нужен поиск длины шага?} Вдали от минимума квадратичная модель может быть неточной. Выбирая $\alpha_k>0$, например по условию Армихо,
\[
f(x^k+\alpha_kp_k)\le f(x^k)+c_1\alpha_k g(x^k)\T p_k,\qquad 0<c_1<1,
\]
обеспечивают убывание функции. Для направления спуска условие выполняется при достаточно малом шаге. При условиях локальной теоремы и $c_1<1/2$ вблизи минимума принимается полный шаг и сохраняется квадратичная скорость.
\end{document}
